\documentclass[tikz,border=6pt]{standalone}
\usepackage{ctex}
\usepackage{amsmath}
\usetikzlibrary{calc, arrows.meta}

\begin{document}
\begin{tikzpicture}[scale=1.1, thick]

% 焦半径示意图：|PF| = x_0 + p/2
% 抛物线 y^2 = 4x (p=2)，焦点 F(1,0)，准线 x=-1

% 坐标轴
\draw[->, thin, gray] (-2.2, 0) -- (5.2, 0) node[right, font=\small] {$x$};
\draw[->, thin, gray] (0, -3.2) -- (0, 3.2) node[above, font=\small] {$y$};
\node[font=\footnotesize, below right] at (0,0) {$O$};

% 抛物线
\draw[black, thick, domain=-2.8:2.8, smooth, samples=80]
  plot ({(\x)^2/4}, \x);

% 准线 x = -1
\draw[red, thick] (-1, -3.0) -- (-1, 3.0);
\node[red, font=\small, left] at (-1, 2.5) {准线 $l$};
\node[red, font=\small, below] at (-1, -3.0) {$x=-\dfrac{p}{2}$};

% 焦点 F
\fill[black] (1, 0) circle [radius=2.5pt];
\node[below, font=\small] at (1, -0.05) {$F\!\left(\dfrac{p}{2},0\right)$};

% 动点 P（参数 t=2.2，P≈(1.21, 2.2)）
\coordinate (P) at ({(2.2)^2/4}, 2.2); % (1.21, 2.2)
\fill[black] (P) circle [radius=2.5pt];
\node[above right, font=\small] at (P) {$P(x_0, y_0)$};

% A 是 P 在准线的射影
\coordinate (A) at (-1, 2.2);
\fill[black] (A) circle [radius=2.5pt];
\node[above left, font=\small] at (A) {$A$};

% 焦半径 PF（蓝色虚线）
\draw[blue, dashed, thick] (P) -- (1, 0)
  node[midway, right, font=\footnotesize] {$|PF|$};

% PA（到准线距离，红色实线）
\draw[red, thick] (P) -- (A)
  node[midway, above, font=\footnotesize] {$|PA|$};

% 直角标记
\draw[thin] (-1+0.14, 2.2) -- (-1+0.14, 2.2-0.14) -- (-1, 2.2-0.14);

% x_0 的水平距离标注
\draw[gray, dashed, thin] (P) -- ({(2.2)^2/4}, 0);
\node[gray, font=\footnotesize, below] at ({(2.2)^2/4}, 0) {$x_0$};

% 虚线从 A 到 x 轴
\draw[gray, dashed, thin] (A) -- (-1, 0);
\node[gray, font=\footnotesize, below] at (-1, -0.05) {$-\dfrac{p}{2}$};

% 焦半径公式标注
\node[font=\small, align=center, draw=black, rounded corners=3pt, fill=white] at (3.2, -2.5)
  {$|PF| = |PA| = x_0 + \dfrac{p}{2}$};

\end{tikzpicture}
\end{document}
